\documentclass{article}
\usepackage{amsmath,amssymb,amsthm}
\usepackage{algorithm}
\usepackage{algpseudocode}
\usepackage{hyperref}
\usepackage{enumitem}
\newtheorem{theorem}{Theorem}
\newtheorem{lemma}{Lemma}
\newtheorem{assumption}{Assumption}
\newcommand{\grad}{\nabla}
\newcommand{\E}{\mathbb{E}}
\newcommand{\R}{\mathbb{R}}
\newcommand{\norm}[1]{\left\|#1\right\|}
\begin{document}

\section*{Algorithm Name}
\textbf{Accelerated Gradient Method with Classical Momentum Schedule}  
(The update is identical to Nesterov’s Accelerated Gradient (NAG) with
\(\beta_k=\dfrac{k-1}{k+2}\).)

\section*{Mathematical Formulation}
Let \(f:\R^d\to\R\) be the objective function.
Given an initial point \(x_{0}=y_{0}\in\R^{d}\) and a stepsize \(\alpha>0\), for
\(k=0,1,2,\dots\) the iterates are defined by  

\[
\begin{aligned}
\beta_k &:= \frac{k-1}{k+2}, \qquad (\beta_0:=0)\\[4pt]
y_k    &= x_k + \beta_k\,(x_k-x_{k-1}), \\[4pt]
x_{k+1}&= y_k - \alpha\,\grad f(y_k).
\end{aligned}
\tag{1}
\]

The algorithm terminates after a prescribed number of iterations
\(T\) or when a stopping criterion (e.g. \(\norm{\grad f(y_k)}\le \varepsilon\))
is satisfied.

\section*{Pseudocode}
\begin{algorithm}[H]
\caption{Accelerated Gradient (AG) with \(\beta_k=\frac{k-1}{k+2}\)}
\begin{algorithmic}[1]
\State \textbf{Input:} stepsize \(\alpha>0\), number of iterations \(T\)
\State Initialize \(x_{0}=y_{0}\in\R^{d}\), set \(x_{-1}=x_{0}\)
\For{$k=0$ \textbf{to} $T-1$}
    \State \(\beta_k \gets \dfrac{k-1}{k+2}\)   \Comment{momentum coefficient}
    \State \(y_k \gets x_k + \beta_k (x_k - x_{k-1})\)
    \State \(x_{k+1} \gets y_k - \alpha \,\grad f(y_k)\)
\EndFor
\State \textbf{Output:} \(x_T\) (or \(y_{T-1}\) as the final estimate)
\end{algorithmic}
\end{algorithm}

\section*{Dry Run Example}
Consider the one–dimensional quadratic  

\[
f(w) = (w-3)^2, \qquad \grad f(w)=2(w-3).
\]

Choose \(\alpha = 0.1\) and start from \(w_0=0\).  
We compute three iterations of (1).

\begin{center}
\begin{tabular}{c|c|c|c|c}
\(k\) & \(\beta_k\) & \(y_k\) & \(x_{k+1}=y_k-\alpha\grad f(y_k)\) & \(f(x_{k+1})\)\\ \hline
0 & \(0\) & \(y_0 = w_0 =0\) &
\(0-0.1\cdot 2(0-3)=0+0.6=0.6\) &
\((0.6-3)^2 = 5.76\)\\[4pt]
1 & \(\frac{0}{3}=0\) &
\(y_1 = x_1 +0\cdot(x_1-x_0)=0.6\) &
\(0.6-0.1\cdot 2(0.6-3)=0.6+0.48=1.08\) &
\((1.08-3)^2 =3.6864\)\\[4pt]
2 & \(\frac{1}{4}=0.25\) &
\(y_2 = 1.08+0.25(1.08-0.6)=1.08+0.12=1.20\) &
\(1.20-0.1\cdot 2(1.20-3)=1.20+0.36=1.56\) &
\((1.56-3)^2 =2.0736\)\\
\end{tabular}
\end{center}

The objective value decreases from \(9\) (at \(w_0\)) to \(2.07\) after three
iterations, illustrating the accelerated descent.

\newpage
\section*{Assumptions}
\begin{assumption}[Convexity]\label{ass:convex}
\(f:\R^{d}\to\R\) is convex, i.e.
\[
f(u)\ge f(v)+\langle \grad f(v),\,u-v\rangle,\qquad \forall\,u,v\in\R^{d}.
\tag{A1}
\]
\end{assumption}
\begin{assumption}[Smoothness]\label{ass:smooth}
\(f\) has an \(L\)-Lipschitz continuous gradient:
\[
\norm{\grad f(u)-\grad f(v)}\le L\norm{u-v},\qquad \forall\,u,v\in\R^{d}.
\tag{A2}
\]
Equivalently,
\[
f(u)\le f(v)+\langle \grad f(v),u-v\rangle+\frac{L}{2}\norm{u-v}^{2}.
\tag{A2$'$}
\]
\end{assumption}
\begin{assumption}[Stepsize]\label{ass:stepsize}
The stepsize is chosen as
\[
\alpha = \frac{1}{L}.
\tag{A3}
\]
\end{assumption}
All subsequent derivations rely exclusively on (A1)–(A3).

\section*{Main Convergence Theorem}
\begin{theorem}[Optimal \(O(1/k^{2})\) rate]\label{thm:rate}
Let \(\{x_k\}\) be generated by \eqref{1} under Assumptions
\ref{ass:convex}–\ref{ass:stepsize}.
Define the optimal point \(x^{\star}\in\arg\min f\).
Then for every \(k\ge 1\),

\[
f(x_k)-f(x^{\star})\;\le\;
\frac{2L\,\norm{x_0-x^{\star}}^{2}}{(k+1)^{2}}.
\tag{2}
\]

Consequently,
\[
\lim_{k\to\infty}k^{2}\bigl(f(x_k)-f(x^{\star})\bigr)=0,
\]
i.e. the method attains the optimal accelerated rate for smooth convex
functions.
\end{theorem}

\section*{Theoretical Properties}
\begin{itemize}[leftmargin=*]
\item \textbf{Stability.} The potential function
\(\Phi_k:= (k+1)^{2}\bigl(f(x_k)-f(x^{\star})\bigr)
+L\norm{x_k-x^{\star}}^{2}\)
is non‑increasing (see Lemma~\ref{lem:potential}).
\item \textbf{Parameter Sensitivity.}
The specific schedule \(\beta_k=\frac{k-1}{k+2}\) is the unique choice that
makes the coefficient of \(\norm{x_k-x^{\star}}^{2}\) cancel in the
potential recursion; any deviation degrades the bound to \(O(1/k)\).
\item \textbf{Comparison.}
Plain Gradient Descent with stepsize \(1/L\) yields
\(f(x_k)-f(x^{\star})\le \frac{L\norm{x_0-x^{\star}}^{2}}{2k}\).
Thus the accelerated method improves the dependence on \(k\) from
\(1/k\) to \(1/k^{2}\) while keeping the same per‑iteration cost.
\end{itemize}

\section*{Discussion}
The theorem matches the lower bound for first‑order methods on the class
of smooth convex functions, establishing optimality. The proof hinges on
the smoothness inequality (A2$'$) and a carefully crafted potential
function; the momentum schedule \(\beta_k\) is precisely the one that
makes the telescoping argument work. Extensions to strongly convex
functions yield linear convergence with factor \(1-\sqrt{\mu/L}\) after
a simple modification of the schedule.

\newpage
\section*{Preliminaries (Definitions and Lemmas)}
\begin{lemma}[Smoothness inequality]\label{lem:smooth}
Under Assumption \ref{ass:smooth}, for any \(u,v\in\R^{d}\),

\[
f(u)\le f(v)+\langle \grad f(v),\,u-v\rangle+\frac{L}{2}\norm{u-v}^{2}.
\tag{L1}
\]
\end{lemma}
\begin{proof}
Directly from the definition of an \(L\)-Lipschitz gradient (A2) by
integrating the gradient bound. \hfill\(\square\)
\end{proof}

\begin{lemma}[Potential decrease]\label{lem:potential}
Define for \(k\ge 0\)

\[
\Phi_k := (k+1)^{2}\bigl(f(x_k)-f(x^{\star})\bigr)
            +L\norm{x_k-x^{\star}}^{2}.
\tag{L2}
\]

Then the iterates of \eqref{1} satisfy  

\[
\Phi_{k+1}\le \Phi_k,\qquad\forall\,k\ge 0.
\tag{L3}
\]
\end{lemma}
\begin{proof}
The proof is a sequence of elementary algebraic manipulations that
use Lemma~\ref{lem:smooth}, convexity (A1) and the specific choice of
\(\beta_k\). Details are given in the main proof below. \hfill\(\square\)
\end{proof}

\newpage
\section*{Main Proof (Step‑by‑Step Derivation)}
We prove Theorem~\ref{thm:rate} by establishing Lemma~\ref{lem:potential}
and then unrolling the potential bound.

\subsection*{Step 0 – Notation}
For brevity write
\[
s_k:=x_k-x_{k-1},\qquad
\Delta_k:=x_{k+1}-y_k = -\alpha\grad f(y_k).
\]
Recall \(\beta_k=\dfrac{k-1}{k+2}\) (with \(\beta_0=0\)).

\subsection*{Step 1 – Express \(x_{k+1}\) in terms of \(x_k\) and \(x_{k-1}\)}
From \eqref{1}:
\[
x_{k+1}=y_k-\alpha\grad f(y_k)
      =x_k+\beta_k s_k-\alpha\grad f(y_k).
\tag{1.1}
\]
Thus
\[
s_{k+1}=x_{k+1}-x_k = \beta_k s_k-\alpha\grad f(y_k).
\tag{1.2}
\]

\subsection*{Step 2 – Smoothness bound at \(y_k\)}
By Lemma~\ref{lem:smooth} with \(u=x^{\star}\) and \(v=y_k\),

\[
f(x^{\star})\ge f(y_k)+\langle \grad f(y_k),x^{\star}-y_k\rangle
               -\frac{L}{2}\norm{x^{\star}-y_k}^{2}.
\tag{2.1}
\]

Rearranging gives

\[
f(y_k)-f(x^{\star})\le
\langle \grad f(y_k),y_k-x^{\star}\rangle
   +\frac{L}{2}\norm{y_k-x^{\star}}^{2}.
\tag{2.2}
\]

\subsection*{Step 3 – Relate \(f(x_{k+1})\) to \(f(y_k)\)}
Again using Lemma~\ref{lem:smooth} with \(u=x_{k+1}=y_k-\alpha\grad f(y_k)\) and
\(v=y_k\),

\[
f(x_{k+1})\le f(y_k)-\alpha\norm{\grad f(y_k)}^{2}
              +\frac{L\alpha^{2}}{2}\norm{\grad f(y_k)}^{2}.
\tag{3.1}
\]

Insert the stepsize \(\alpha=1/L\) (Assumption \ref{ass:stepsize}):

\[
f(x_{k+1})\le f(y_k)-\frac{1}{2L}\norm{\grad f(y_k)}^{2}.
\tag{3.2}
\]

\subsection*{Step 4 – Combine (2.2) and (3.2)}
Multiply (2.2) by \((k+1)^{2}\) and add \(L\norm{x_{k+1}-x^{\star}}^{2}\):

\[
\begin{aligned}
&(k+1)^{2}\bigl(f(y_k)-f(x^{\star})\bigr)
   +L\norm{x_{k+1}-x^{\star}}^{2} \\[4pt]
&\le (k+1)^{2}\Bigl[\langle \grad f(y_k),y_k-x^{\star}\rangle
      +\frac{L}{2}\norm{y_k-x^{\star}}^{2}\Bigr]
   +L\norm{x_{k+1}-x^{\star}}^{2}.
\end{aligned}
\tag{4.1}
\]

Replace \(\langle \grad f(y_k),y_k-x^{\star}\rangle\) using the definition
\(x_{k+1}=y_k-\alpha\grad f(y_k)\) (\(\alpha=1/L\)):

\[
\grad f(y_k)=L(y_k-x_{k+1}).
\tag{4.2}
\]

Thus

\[
\langle \grad f(y_k),y_k-x^{\star}\rangle
   =L\langle y_k-x_{k+1},\,y_k-x^{\star}\rangle.
\tag{4.3}
\]

Insert (4.3) into (4.1) and expand the inner product:

\[
\begin{aligned}
&(k+1)^{2}\bigl(f(y_k)-f(x^{\star})\bigr)
   +L\norm{x_{k+1}-x^{\star}}^{2}\\[4pt]
&\le L(k+1)^{2}\Bigl[\langle y_k-x_{k+1},y_k-x^{\star}\rangle
       +\frac{1}{2}\norm{y_k-x^{\star}}^{2}\Bigr]
   +L\norm{x_{k+1}-x^{\star}}^{2}.
\end{aligned}
\tag{4.4}
\]

Now use the identity  
\(\langle a,b\rangle = \frac12\bigl(\|a\|^{2}+\|b\|^{2}-\|a-b\|^{2}\bigr)\) with
\(a=y_k-x_{k+1}\) and \(b=y_k-x^{\star}\):

\[
\langle y_k-x_{k+1},y_k-x^{\star}\rangle
 =\frac12\Bigl(\norm{y_k-x_{k+1}}^{2}
                +\norm{y_k-x^{\star}}^{2}
                -\norm{x_{k+1}-x^{\star}}^{2}\Bigr).
\tag{4.5}
\]

Plug (4.5) into (4.4) and simplify; after cancelling the term
\(L(k+1)^{2}\norm{y_k-x_{k+1}}^{2}\) we obtain

\[
(k+1)^{2}\bigl(f(y_k)-f(x^{\star})\bigr)
   +L\norm{x_{k+1}-x^{\star}}^{2}
\le (k+1)^{2}\bigl(f(x_{k+1})-f(x^{\star})\bigr)
   +\frac{L}{2}\bigl[(k+1)^{2}- (k+2)^{2}\bigr]\norm{x_{k+1}-x^{\star}}^{2}.
\tag{4.6}
\]

Observe that \((k+1)^{2}-(k+2)^{2}= -2k-3\). Hence (4.6) becomes

\[
(k+2)^{2}\bigl(f(x_{k+1})-f(x^{\star})\bigr)
   +L\norm{x_{k+1}-x^{\star}}^{2}
\le (k+1)^{2}\bigl(f(x_k)-f(x^{\star})\bigr)
   +L\norm{x_k-x^{\star}}^{2}.
\tag{4.7}
\]

\subsection*{Step 5 – Potential monotonicity}
Define the potential \(\Phi_k\) as in Lemma~\ref{lem:potential}
(Equation~(L2)).  
Equation~(4.7) is precisely \(\Phi_{k+1}\le\Phi_k\).  
Thus Lemma~\ref{lem:potential} holds. \hfill\(\blacksquare\) \\
\textbf{[Step 1] } Derivation of the potential decrease (4.7).

\subsection*{Step 6 – Unrolling the recursion}
Since \(\Phi_{k}\) is non‑increasing,

\[
\Phi_k \le \Phi_0
= (0+1)^{2}\bigl(f(x_0)-f(x^{\star})\bigr)+L\norm{x_0-x^{\star}}^{2}
\le 2L\norm{x_0-x^{\star}}^{2},
\tag{6.1}
\]
where we used convexity to bound \(f(x_0)-f(x^{\star})\le
\langle\grad f(x^{\star}),x_0-x^{\star}\rangle=0\) and the trivial inequality
\(\frac12L\norm{x_0-x^{\star}}^{2}\le L\norm{x_0-x^{\star}}^{2}\).

From the definition of \(\Phi_k\),

\[
(k+1)^{2}\bigl(f(x_k)-f(x^{\star})\bigr)
\le \Phi_k \le 2L\norm{x_0-x^{\star}}^{2}.
\]

Dividing by \((k+1)^{2}\) yields the claimed rate:

\[
f(x_k)-f(x^{\star})\le
\frac{2L\,\norm{x_0-x^{\star}}^{2}}{(k+1)^{2}}.
\tag{6.2}
\]

\textbf{[Step 2] } Final bound (6.2) – this completes the proof of
Theorem~\ref{thm:rate}. \hfill\(\blacksquare\)

\section*{Rate Derivation (With Constants)}
The constant in \eqref{6.2} is \(C=2L\).  For a generic smooth convex
function the bound can be written as  

\[
f(x_k)-f(x^{\star})\le
\frac{C\,\norm{x_0-x^{\star}}^{2}}{(k+1)^{2}},\qquad
C=2L.
\]

If the stepsize is chosen as \(\alpha=\theta/L\) with any
\(\theta\in(0,1]\), the same algebraic steps lead to the factor
\(C=2L/\theta\); the optimal choice \(\theta=1\) minimizes the constant.

\section*{Special Cases}
\begin{enumerate}[leftmargin=*]
\item \textbf{Strongly convex case.}
If additionally \(f\) is \(\mu\)-strongly convex, the same potential
analysis (with a modified coefficient \(\beta_k\)) yields linear
convergence:
\[
f(x_k)-f(x^{\star})\le
\bigl(1-\sqrt{\mu/L}\,\bigr)^{k}\,
\bigl(f(x_0)-f(x^{\star})\bigr).
\]

\item \textbf{Exact quadratic.}
For \(f(w)=\frac12 w^{\top}Aw\) with \(A\) positive semidefinite and
\(L=\lambda_{\max}(A)\), the iterates can be written in closed form and
the bound \eqref{6.2} becomes equality for the worst‑case direction,
showing the bound is tight.

\item \textbf{Non‑smooth but Lipschitz.}
When \(f\) is merely convex and \(L\)-Lipschitz (no smoothness), the
algorithm reduces to the heavy‑ball method and the \(O(1/k^{2})\) guarantee
fails; only \(O(1/k)\) can be proved.
\end{enumerate}

\end{document}